\chapter{TUJUAN DAN MANFAAT PENELITIAN}
\label{bab:3}

\section{Tujuan Penelitian}
Berdasarkan rumusan masalah yang telah diuraikan sebelumnya, tujuan utama yang ingin dicapai dari pelaksanaan penelitian tugas akhir ini adalah sebagai berikut:
\begin{enumerate}
    \item Merancang, mensimulasikan, dan menguji algoritma pemrograman kinematika gerak berorde tiga untuk menggenerasikan profil pergerakan \textit{trapezoidal} dan \textit{S-Curve} secara bergantian pada unit pemroses mikrokontroler \textit{dual-core} secara \textit{real-time}.
    \item Membuat desain dan mengimplementasikan modul perangkat keras tambahan berupa rangkaian elektronik \textit{differential line driver} berbasis chip transceiver industri guna mengonversi pulsa kontrol menjadi sinyal diferensial berstandar RS-422 menuju \textit{Servo Amplifier} Melservo-J2-Super.
    \item Menganalisis parameter performa dinamis motor servo AC secara komparatif berdasarkan data empiris nilai kesalahan posisi (\textit{positional error}), durasi \textit{settling time} transien, serta memvalidasi kehalusan pergerakan meredam sentakan mekanis melalui visualisasi metode uji \textit{liquid slosh control} menggunakan media cairan pada modul ED4270m.
\end{enumerate}

\section{Manfaat Penelitian}
Pelaksanaan riset rancang bangun dan analisis komparatif ini diharapkan mampu memberikan kontribusi yang signifikan, baik dari aspek pengembangan keilmuan maupun aplikasi praktis di lapangan:
\begin{enumerate}
    \item \textbf{Manfaat Teoritis (Akademik):}
    \begin{itemize}
        \item Memberikan kontribusi referensi ilmiah mengenai pemodelan matematika terapan untuk pembatasan nilai \textit{jerk} pada sistem kendali posisi non-linier menggunakan perangkat keras berbiaya rendah.
        \item Menyediakan paradigma baru dalam metode pemrograman sistem tertanam mengenai teknik distribusi beban kerja komputasi numerik berat memanfaatkan teknologi \textit{multitasking dual-core} secara deterministik.
    \end{itemize}
    \item \textbf{Manfaat Praktis (Aplikatif):}
    \begin{itemize}
        \item Menyediakan alternatif sistem kendali gerak mandiri berbiaya murah namun memiliki keandalan tinggi dan kebal terhadap derau elektro-mekanis setara standar industri komersial.
        \item Membantu memperpanjang usia pakai dari komponen mekanis presisi tinggi pada industri, karena keberhasilan algoritma \textit{S-Curve} secara langsung meminimalkan keausan akibat hantaman gaya inersia getaran balik.
        \item Menghasilkan purwarupa fisik sistem pengontrol beserta draf dokumen analitis komparatif yang siap diintegrasikan sebagai modul praktikum kendali posisi presisi tinggi di Laboratorium Mekatronika.
    \end{itemize}
\end{enumerate}